\section{Easy、Medium、Hard：为什么单次解码会低估分布能力}

讲者用三道 HumanEval problem 展示 pass@1 的数量级差异。Slide 的价值不在于记住具体实现，而在于观察 prompt ambiguity、algorithm depth 与 rare correctness 如何改变 sampling strategy。

\subsection{Easy problem}

本节先看 easy case，建立“正确样本位于高概率区域”时的基线，再逐步增加 algorithm 和 boundary 难度。读代码时先根据 docstring 手工列出 specification，再观察模型是否使用常见模板就能覆盖 normal 与 edge cases。

\lecturefigure{slide-34-easy-problem.jpg}{Easy problem 的正确实现位于模型分布高概率区域。}{官方视频 38:03。}

\begin{knowledgebox}{读图：高 pass@1 仍需读 specification}
代码短、pattern 常见、boundary 少时，greedy 或低温 sample 容易正确。先检查 prompt 是否完整，再检查 candidate 是否处理 empty/negative/duplicate。高分可能来自训练中常见模板，不代表模型理解所有变体。
\end{knowledgebox}

\subsection{Medium problem}

接下来观察 medium case：主体结构可能正确，但一个 branch 或 off-by-one 就足以让全部 hidden tests 失败。这类题最能说明自然语言流畅度与执行正确性脱钩，也最适合分析多次 sampling 是否产生互补错误。

\lecturefigure{slide-35-medium-problem.jpg}{Medium problem 要求更多步骤或边界处理，正确样本概率明显下降。}{官方视频 38:36。}

\begin{knowledgebox}{读图：错误可能集中在一个 branch}
候选通常已有正确框架，却在循环范围、条件顺序或 state update 上失败。此时多次 sampling 可能探索到正确变体；也可用 test feedback、static analysis 或 repair loop 提高选择。Benchmark 只统计最终 pass，不解释错误类型，需要额外 taxonomy。
\end{knowledgebox}

\subsection{Hard problem}

进一步进入 hard case，问题不再是“默认输出是否正确”，而是正确 program 是否在分布中拥有足够 probability mass 供 search 找到。若多数 samples 共享同一错误 algorithm，增加 $k$ 只会重复失败；若错误多样，verifier 才可能从长尾中筛出正确候选。

\lecturefigure{slide-36-hard-problem.jpg}{Hard problem 的单样本正确率很低，但分布中仍可能存在正确实现。}{官方视频 39:00。}

\begin{knowledgebox}{读图：低 pass@1 与零能力不同}
若 pass@1 约为千分量级，单次 demo 几乎总失败；大量 samples 可能仍发现正确程序。系统必须决定能否承担生成、执行与筛选成本。若没有可靠 verifier，large-k 只会产生更多不可审计代码。
\end{knowledgebox}

\subsection{为什么需要 sampling-aware metric}

下面从单题例子转向全 benchmark 分布。直方图展示有些题几乎总会、有些存在长尾；平均 pass@1 会隐藏“少量候选即可解决”与“再多样本也不行”的区别。这里要携带的问题是：额外 generation budget 应投给哪些题，何时应改进模型而不是继续采样。

\lecturefigure{slide-37-need-pass-at-k-hist.jpg}{题目级成功概率呈长尾，单一平均分隐藏可采样性差异。}{官方视频 40:30。}

\begin{knowledgebox}{读图：横轴题目成功率，纵轴题目数量}
若大量题集中在接近 0 或 1，模型表现呈 bimodal；中间区域最适合通过 sampling 提升。需要报告 per-problem distribution，而非只给 mean。还应检查不同 model size 的 mass 如何移动，以及难题是否真正获得非零概率。
\end{knowledgebox}

\lecturefigure{slide-38-need-pass-at-k-scatter.jpg}{不同题目的 pass@1 与 pass@k 提升并不均匀。}{官方视频 40:57。}

\begin{knowledgebox}{读图：散点偏离对角线意味着 sampling value}
若点在高-k 轴显著上移，说明多样采样可找到 rare correct programs；仍贴近零的点表示 distribution 几乎没有正确质量。Sampling 不是万能，它只能放大已有 probability mass，不能凭空创造模型从未学会的 algorithm。
\end{knowledgebox}

\subsection{本章小结}

题目难度改变正确样本在分布中的位置。Pass@1 测默认一次命中，pass@k 测生成与筛选系统的搜索潜力；两者都需要 per-problem 解释。

